\section{The \texorpdfstring{\SxC}{SxC} Decomposition}
\label{sec:method}

\paragraph{Setup.} ROME inserts a key$\rightarrow$value association by a rank-one
update to a single MLP down-projection $W$. Given an edit key $k$ (the
subject's last-token key at the critical layer) and target value $v$, the
update is
\begin{equation}
\dW = \frac{(v-Wk)\,k^{\top}}{k^{\top}k}
    = \frac{\mathrm{outer}(v-Wk,\,k)}{\|k\|^2},
\label{eq:rome}
\end{equation}
the minimum-norm solution to $(W+\dW)k=v$.

\paragraph{Collateral perturbation to an unrelated probe.} For an unrelated
fact whose key at the same layer is $k_p$, the edit perturbs its value
read-out by $\dW\cdot k_p = (v-Wk)\,(k^{\top}k_p)/\|k\|^2$. Taking magnitudes
and writing $k^{\top}k_p = \|k\|\,\|k_p\|\cos(k,k_p)$,
\begin{equation}
\|\dW\cdot k_p\| = \underbrace{\frac{\|v-Wk\|}{\|k\|}}_{S}\;\|k_p\|\;
                   \underbrace{|\cos(k,k_p)|}_{|C|},
\label{eq:sc}
\end{equation}
where $S$ is a per-\emph{edit} strength (constant across probes), $|C|$ is the
per-(edit,probe) key-geometry, and $\|k_p\|$ is a per-\emph{probe} scalar. The
collateral perturbation factorizes multiplicatively into edit strength, probe
norm, and the edit$\times$probe key-cosine. The observed \emph{damage} -- the
drop in the probe's correct-token logit -- is a monotone function of
$\|\dW\cdot k_p\|$ under local linearization of the head, so we measure the
\emph{rank} relationship (a within-probe Spearman) rather than assert exact
proportionality.

\paragraph{\texorpdfstring{\SxC}{SxC} is the closed-form reduction of
first-order gradient influence.} Up to the per-probe norm, $S\cdot|C|$ is the
exact rank-one reduction of a first-order (GradSim/TracIn-style)
training-influence measure \citep{pruthi2020gradsim} of an edit on a probe --
obtainable algebraically, with no backprop (a full derivation and its
assumptions are deferred to a dedicated theorem pass). At L8 the within-probe
$\rho$ for \SxC{} is \scRhoLeight{}, at L10 \scRhoLten{}, at L12
\scRhoLtwelve{}, and at L14 \scRhoLfourteen{} -- versus raw key-cosine
\gWithinLeight{}, \gWithinLten{}, \gWithinLtwelve{}, \gWithinLfourteen{},
respectively. \textbf{\SxC{} beats raw key-cosine at all four layers} once $S$
is normalized per Eq.~\ref{eq:sc} ($S=\|v-Wk\|/\|k\|$, not the unnormalized
residual norm alone); the margin is largest where $S$ carries additional
variance (L12/L14). We frame \SxC{} as the mechanistic, zero-cost surrogate
that explains why cosine works and predicts where its predictive power
changes -- not as a quantity independently validated against a
separately-computed backprop gradient similarity: the on-disk ``GradSim''
baseline (Figure~\ref{fig:surrogate}) is built from the same closed-form
factorization used to define \SxC{} itself, so its agreement with \SxC{} is a
restatement of the same construction rather than an independent empirical
check; a true-backprop GradSim comparison is left to future work.

\subsection{Confound-clean measurement}
\label{sec:gate}
The empirical test of Eq.~\ref{eq:sc} is a \emph{within-probe} partialled
Spearman between $|C|$ and damage, computed down each probe column (holding the
probe fixed differences out $\|k_p\|$ and the probe's intrinsic
susceptibility). We run ROME on Llama-3.2-1B \citep{grattafiori2024llama3},
CounterFact, $N{=}200$ edits
$\times$ $M{=}500$ probes, with two confound controls: the \emph{known-fact
filter} (restrict to probes the base model assigns the true answer,
$\text{pre\_p}>0.05$) and the \emph{successful-edit filter} (drop failed edits). This retires two objections the design pre-registered as
mandatory: (i) \emph{probe-marginal leakage} -- the pooled $\rho$ ($\sim$0.40
at L8) barely drops to the within-probe $\rho$ (\gWithinLeight{}), so the
signal is edit-specific pairwise geometry, not ``which probe is fragile''; and
(ii) \emph{non-independence} -- a column-permutation null (shuffle within
probe) gives permutation-$p$ at the \gPermFloor{} floor. The signal survives
row-marginal double-centering and survives partialling norm-growth out per
column. Full gate results are in \S\ref{sec:regime} and
Figure~\ref{fig:layerlaw}.
